\section*{Main statement and conditional consequences}
\setcounter{theorem}{2}
\begin{theorem}[Main Theorem]
\label{thm:main-intro}
Assume that $\lambda= (\lambda_1,\ldots,\lambda_\ell)$ is such that $\lambda_i \geq \lambda_{i+1}+(\ell-1)$ for all $i \in [\ell-1]$, also $\lambda_{\ell} \geq \ell$ and let $d = \sum_i \lambda_i$. Then, for $n \geq \lambda_1 + \ell$, if $s_\lambda(x_1,\ldots,x_n)$ has an algebraic formula of size $s$ and depth $\Delta$, then the $\ell\times \ell$ determinant $(\mathrm{det}_\ell )$ has an algebraic formula of size $\poly(s)$ and depth $\Delta + O(1).$
\end{theorem}

For suitable choices of $\ell,d,n$ above, we can ensure that these parameters are all polynomially related. The theorem then implies that the Schur polynomials do not have algebraic formulas of polynomial size unless the entire complexity class $\mathrm{VBP}$ collapses to the complexity class $\mathrm{VF}$ which consists of polynomials with small formulas. Moreover, the Schur polynomials do not have  \emph{constant-depth} formulas of subexponential size unless the determinant does, a result that would greatly improve the state-of-the-art in this direction~\cite{r16}. 

The above theorem and its proof have several interesting aspects that we now elaborate on. 

\begin{lemma}[Main Technical Lemma (informal)]\label{lem:key-intro}
Let $g_1,\ldots,g_n\in \mathbb{C}[x_1,\ldots,x_n]$ be ``well-behaved'' algebraically independent polynomials. Then, for any homogeneous polynomial $\tilde{f}$, if $f = \tilde{f}(g_1,\ldots,g_n)$ has a formula of size $s$ and depth $\Delta$, the polynomial $\tilde{f}$ has a formula of size $\poly(s)$ and depth $\Delta + O(1).$
\end{lemma}
\paragraph*{Generalized Vandermonde determinants.} The Vandermonde matrix $(x_i^{n-j})_{i,j\in [n]}$ and its determinant are ubiquitous in computation because of their relation to polynomial interpolation. More precisely, the problem of finding a degree-$(n-1)$ univariate polynomial that takes prescribed values at a specified set of $n$ distinct points involves solving a linear system of equations where the underlying matrix is precisely the Vandermonde matrix. It is, therefore, an important fact that the Vandermonde determinant is computationally much easier than the general determinant: in fact, it has the following standard depth-$2$ formula
\[
\mathrm{det}_n\left((x_i^{n-j})_{i,j\in [n]}\right) = \prod_{i,j\in [n]: i < j}(x_i - x_j).
\]
However, it is unclear whether such small expressions continue to exist if we allow the exponents of the variables to vary more generally. For integers $\mu_1 > \mu_2 > \cdots > \mu_n \geq 0$, consider the \emph{generalized Vandermonde matrix} $(x_i^{\mu_j})_{i,j\in [n]}$. Similar to the Vandermonde matrix, the determinant of this matrix is related to the problem of \emph{sparse polynomial interpolation}, where we are trying to interpolate a polynomial only involving  the monomials of degree $\mu_1,\ldots,\mu_n$ through the given points. 

Can we expect that computing any generalized Vandermonde determinant is much easier than computing the determinant itself? It follows from Theorem~\ref{thm:main-intro} and the bialternant formula from Definition~\ref{def: Schur defn} above that the answer to this question is negative: for certain (polynomially large) exponents, the generalized Vandermonde determinant is not much easier than the determinant.




\setcounter{section}{1}
